\begin{definition}\label{def:2.1.4}
$\Sigma$ 中 $\tau$ 的不动点称为 $(\Sigma,\tau)$ 的\emph{实点}，$\Sigma^\tau$ 的连通分支
称为\emph{实点圆周}。
\end{definition}

走向实曲线拓扑分类的第一步由 Harnack 不等式给出\footnote{事实上，Harnack 只对实平面
曲线证明了定理 2.1.5，而任意实曲线的情形是克莱因证明的，这一点由 Erwan Brugallé
告知作者。}，它对 $\Sigma^\tau$ 的连通分支个数给出拓扑上界（所谓拓扑，是指它不依赖
于 $\Sigma$ 上复结构的选择，只依赖于亏格）。我们给出一个微分几何的证明，其中用到
曲面的分类（例如可参看 [4]）。

\begin{theorem}[Harnack 不等式，[7]]\label{thm:2.1.5}
设 $(\Sigma,\tau)$ 是 Real 黎曼面，$\Sigma^\tau$ 有 $k$ 个连通分支。则 $k\le g+1$。
若 $k=g+1$，则 $\Sigma/\tau$ 同胚于球面挖去 $g+1$ 个圆盘。
\end{theorem}

应当指出，定理 2.1.5 的完整表述还断言上述界是最优的：对任意 $g\ge0$，存在 Real
黎曼面 $(\Sigma,\tau)$ 使 $\Sigma^\tau$ 恰有 $g+1$ 个连通分支。

\begin{lemma}\label{lem:2.1.6}
$\Sigma$ 与 $\Sigma/\tau$ 的欧拉示性数满足
\[
\chi(\Sigma/\tau)=\tfrac12\chi(\Sigma)=1-g .
\]
\end{lemma}

\begin{proof}
典范投影 $p:\Sigma\to\Sigma/\tau$ 在 $\partial(\Sigma/\tau)$ 之外是二对一的，并且限制
在 $\Sigma^\tau$ 上给出微分同胚 $\Sigma^\tau\simeq\partial(\Sigma/\tau)$。从 $\Sigma/\tau$
的一个三角剖分 $T$ 出发，不与 $\partial(\Sigma/\tau)$ 相交的任一三角形（不妨称之为
(a) 型）在 $p$ 下的拉回给出 $\Sigma$ 中两个不相交的三角形（只要原三角剖分 $T$ 中的
三角形足够小，这总可假定）。有一个顶点落在 $\partial(\Sigma/\tau)$ 上的三角形（称之为
(b) 型）的拉回给出 $\Sigma$ 中恰有一个公共顶点的两个三角形。而有一条边落在
$\partial(\Sigma/\tau)$ 上的三角形（称之为 (c) 型）的拉回给出 $\Sigma$ 中恰有一条公共
边的两个三角形。由于 $\partial(\Sigma/\tau)$ 是 $k$ 个圆周之并，其欧拉示性数为零。故
在三角剖分 $T$ 中，(b) 型三角形的个数与 (c) 型相同。特别地，$T$ 中落在
$\partial(\Sigma/\tau)$ 内的顶点数 $v_1$ 等于落在其内的边数 $e_1$。记 $T$ 中其余的
顶点数（相应地，边数、面数）为 $v_2$（相应地，$e_2$、$c$）。则
\[
\chi(\Sigma/\tau)=(v_1+v_2)-(e_1+e_2)+c=v_2-e_2+c,
\]
而用上述拉回三角剖分 $p^*T$ 计算 $\Sigma$：
\[
\chi(\Sigma)=(v_1+2v_2)-(e_1+2e_2)+2c=2v_2-2e_2+2c=2\chi(\Sigma/\tau).\qedhere
\]
\end{proof}

\begin{proof}[定理 2.1.5 的证明]
紧连通曲面 $\Sigma/\tau$ 有 $k$ 个边界分支。由曲面分类，$\chi(\Sigma/\tau)\le2-k$。
事实上，若 $\Sigma/\tau$ 可定向，则它要么是球面挖去 $k$ 个圆盘（此时
$\chi(\Sigma/\tau)=2-k$），要么是 $\hat g$ 个环面的连通和挖去 $k$ 个圆盘（此时
$\chi(\Sigma/\tau)=2-2\hat g-k<2-k$）。若 $\Sigma/\tau$ 不可定向，则它是 $h\ge1$ 个
实射影平面 $\RP^2$ 的连通和挖去 $k$ 个圆盘（此时 $\chi(\Sigma/\tau)=2-h-k<2-k$）。
于是由引理 2.1.6，
\[
2-k\ \ge\ \chi(\Sigma/\tau)\ =\ 1-g,
\]
故 $k\le g+1$。此外，若 $k=g+1$，则 $\chi(\Sigma/\tau)=2-k$，而证明过程中已经看到，
这只有在 $\Sigma/\tau$ 同胚于球面挖去 $k$ 个圆盘时才发生。
\end{proof}

由实曲线同态的定义可以清楚看出，亏格 $g$ 与 $\Sigma^\tau$ 的连通分支数 $k$ 都是
$(\Sigma,\tau)$ 的拓扑不变量。事实上，若 $f:\Sigma\to\Sigma'$ 是同胚且
$f\circ\tau=\tau'\circ f$，则 $f(\Sigma^\tau)\subset(\Sigma')^{\tau'}$，且 $f$ 在
$\Sigma^\tau$ 与 $(\Sigma')^{\tau'}$ 的连通分支之间诱导双射。现在令
\begin{equation}\label{eq:2.4}
a=
\begin{cases}
0, & \text{若 } \Sigma\setminus\Sigma^\tau \text{ 不连通},\\
1, & \text{若 } \Sigma\setminus\Sigma^\tau \text{ 连通}.
\end{cases}
\end{equation}
特别地，若 $k=0$，则 $a=1$。实结构 $\tau$ 称为\emph{分离的}（dividing），若 $a=0$；
称为\emph{非分离的}，若 $a=1$。文献中有时也见到相反的约定，但我们不知道有什么理由
偏好其中一种。

本小节的其余部分用来证明 $(g,k,a)$ 是实曲线的完备拓扑不变量。证明这一分类结果的主要
工具是任意给定紧连通曲面 $S$ 的\emph{double}（双重）概念。

\begin{definition}\label{def:2.1.7}
设 $S$ 是紧连通曲面。$S$ 的一个 \emph{double} 是一个三元组 $(P,\sigma,q)$，其中 $P$
是无边的可定向曲面，$\sigma$ 是 $P$ 的反转定向对合，$q:P\to S$ 是诱导同胚
$q:P/\sigma\xrightarrow{\ \simeq\ }S$ 的连续映射。double 的同态是连续映射
$\phi:P\to P'$，使得 $\phi\circ\sigma=\sigma'\circ\phi$。
\end{definition}

注意，反转定向的自微分同胚这一概念在 $P$ 上有意义，即使不选取定向。我们要提醒读者，
定义 2.1.7 的术语并不完全标准（比较 [1, 14]）。不论术语如何，要点是曲面 $P$ 可定向、
无边界，且 $\sigma$ 反转定向。事实上，我们明确希望 $P$ 如下从 $S$ 构造出来。

\medskip
\noindent\emph{第一种情形：$S$ 可定向。}此时 $S$ 同胚于亏格 $\hat g=0$ 的曲面挖去
$r$ 个圆盘（$r\ge0$），或者同胚于 $\hat g$ 个环面的连通和挖去 $r$ 个圆盘
（$\hat g\ge1$，$r\ge0$）。我们用 $\Sigma_{\hat g,r}$ 记这些曲面。取 $\Sigma_{\hat g,r}$
的两个副本 $\Sigma^{(1)}_{\hat g,r}$ 与 $\Sigma^{(2)}_{\hat g,r}$，并令
\[
P:=\Bigl(\Sigma^{(1)}_{\hat g,r}\sqcup\Sigma^{(2)}_{\hat g,r}\Bigr)\Big/\sim,
\]
其中 $x\in\partial\Sigma^{(1)}_{\hat g,r}$ 与 $x\in\partial\Sigma^{(2)}_{\hat g,r}$ 相粘。
特别地，若 $r\ge1$，则 $P$ 是亏格 $2\hat g+r-1$ 的闭可定向曲面。对合 $\sigma$ 由把
$x\in\Sigma^{(1)}_{\hat g,r}$ 送到 $x\in\Sigma^{(2)}_{\hat g,r}$ 定义。由于粘合的方式，
$\sigma$ 反转定向。$\sigma$ 在 $P$ 中的不动点集有 $r$ 个连通分支，而映射
$q:P\to\Sigma_{\hat g,r}$（把 $x\in\Sigma^{(i)}_{\hat g,r}$ 送到 $x\in\Sigma_{\hat g,r}$）
是良定义的，并诱导同胚 $P/\sigma\simeq\Sigma_{\hat g,r}$。注意 $P$ 连通当且仅当
$r\ge1$，即原可定向曲面 $S$ 的边界非空。这正是我们所关心的情形。

\medskip
\noindent\emph{第二种情形：$S$ 不可定向。}此时 $S$ 包含一个开莫比乌斯带，且存在唯一
的整数 $m=1,2$ 以及 $m$ 个互不相交的、含于 $S$ 的开莫比乌斯带 $U_1,\dots,U_m$，使得
紧连通曲面
\[
S':=S\setminus(U_1\cup\cdots\cup U_m)
\]
可定向；见 [4, 11]。若 $S$ 有 $l\ge0$ 个边界分支，则 $S'\simeq\Sigma_{\hat g,l+m}$
对某个 $\hat g\ge0$ 成立。我们记
\[
\partial\Sigma_{\hat g,l+m}=A_1\sqcup\cdots\sqcup A_l\sqcup B_1\sqcup\cdots\sqcup B_m .
\]
每个 $B_j$ 是圆周，因而带有无不动点的对合 $\sigma_j$（想想 $S^1$ 上的 $z\mapsto-z$）。
我们构造
\[
P:=\Bigl(\Sigma^{(1)}_{\hat g,l+m}\sqcup\Sigma^{(2)}_{\hat g,l+m}\Bigr)\Big/\sim,
\]
其中把 $A_i^{(1)}$ 中的点 $x$ 与 $A_i^{(2)}$ 中的 $x$ 等同，并把 $B_j^{(1)}$ 中的点
$x$ 与 $B_j^{(2)}$ 中的点 $\sigma_j(x)$ 等同。特别地，$P$ 是亏格
$2\hat g+(l+m)-1$ 的连通闭可定向曲面。反转定向对合 $\sigma$ 由把
$x\in\Sigma^{(1)}_{\hat g,l+m}$ 送到 $x\in\Sigma^{(2)}_{\hat g,l+m}$ 定义。由于定义粘合
时用到的对合 $\sigma_j$，$\sigma$ 的不动点集有 $l$（而不是 $l+m$）个连通分支（即那些
圆周 $A_i$），且 $\sigma$ 限制在圆周 $B_j$ 上就是 $\sigma_j$。特别地，$P/\sigma$ 同胚于
$S$（而不是 $S'$）。

\medskip
给定紧连通曲面 $S$，我们用 $(P_S,\sigma_S)$ 记如上构造的 $S$ 的 double。注意每个
$(P_S,\sigma_S)$ 联系着一个三元组 $(g_S,k_S,a_S)$，其中 $g_S$ 是 $P$ 的亏格，$k_S$ 是
$\Fix(\sigma_S)$ 的连通分支数（即 $\partial S$ 的连通分支数），而 $a_S=0,1$ 按 (2.4)
定义，取决于 $P_S\setminus P_S^{\sigma_S}$ 是否连通。

对给定的 $S$，$(P_S,\sigma_S)$ 的构造有如下直接推论：
\begin{enumerate}[label=(\roman*), leftmargin=2.2em, itemsep=0.15em]
\item $S$ 不可定向当且仅当 $P_S\setminus P_S^{\sigma_S}$ 连通（即 $a_S=1$）；
\item 两个紧连通曲面 $S$ 与 $S'$ 同胚，当且仅当存在同胚
      $\phi:P_S\to P_{S'}$ 使得 $\phi\circ\sigma_S=\sigma_{S'}\circ\phi$；
\item $S$ 同胚于 $S'$，当且仅当 $(g_S,k_S,a_S)=(g_{S'},k_{S'},a_{S'})$。
\end{enumerate}

于是，考虑到曲面的分类，有下述事实：若 $(P,\sigma)$ 由一个闭可定向曲面与一个反转
定向对合组成，且 $S:=P/\sigma$，则存在同胚 $\phi:P_S\to P$ 使得
$\phi\circ\sigma_S=\sigma\circ\phi$。这样我们就几乎证明了下面这个属于 Weichold 的
结果，它给出 Real 黎曼面的拓扑分类。

\begin{theorem}[Weichold, [20]]\label{thm:2.1.8}
设 $(\Sigma,\tau)$ 与 $(\Sigma',\tau')$ 是两个紧连通 Real 黎曼面。则下列条件等价：
\begin{enumerate}[label=(\alph*), leftmargin=2.2em, itemsep=0.15em]
\item 存在同胚 $\phi:\Sigma\to\Sigma'$ 使得 $\phi\circ\tau=\tau'\circ\phi$；
\item $\Sigma/\tau$ 同胚于 $\Sigma'/\tau'$；
\item $(g,k,a)=(g',k',a')$。
\end{enumerate}
\end{theorem}

三元组 $(g,k,a)$ 称为 Real 黎曼面 $(\Sigma,\tau)$ 的\emph{拓扑型}。它满足如下条件：
\begin{enumerate}[label=(\roman*), leftmargin=2.2em, itemsep=0.15em]
\item 若 $a=0$，则 $1\le k\le g+1$ 且 $k\equiv g+1 \pmod 2$；
\item 若 $a=1$，则 $0\le k\le g$。
\end{enumerate}
注意，Weichold 定理的完整表述还断言：对每个满足上述条件的三元组 $(g,k,a)$，确实存在
拓扑型为 $(g,k,a)$ 的实代数曲线。

\begin{proof}[定理 2.1.8 的证明]
只需再证 $(g,k,a)$ 必须满足定理中列出的条件。由 Harnack 定理 2.1.5 有 $0\le k\le g+1$；
而若 $k=0$，则由 $a$ 的定义 (2.4) 得 $a=1$。若 $a=0$，则 $\Sigma/\tau$ 可定向且同胚于
某个 $\Sigma_{\hat g,k}$（$\hat g\ge0$），故 $\chi(\Sigma/\tau)=2-2\hat g-k$。但由
引理 2.1.6，
\[
\chi(\Sigma/\tau)=\tfrac12\chi(\Sigma)=1-g,
\]
故 $k=g+1-2\hat g\equiv g+1\pmod2$。若 $a=1$，则 $\Sigma/\tau$ 不可定向，因而是
$h\ge1$ 个实射影平面 $\RP^2$ 的连通和挖去 $k$ 个圆盘。于是
$\chi(\Sigma/\tau)=2-h-k$，而同样有 $\chi(\Sigma/\tau)=1-g$，故 $k-g=1-h\le0$。
\end{proof}

\subsection{曲面上的双解析结构}
\label{sec:02.1.3}

本小节中，\emph{曲面}指二维拓扑流形（特别地，Hausdorff 且第二可数），可以带边界。
众所周知，紧连通可定向曲面容许一个一维复解析流形结构（若曲面带边界，该断言应理解为
它是某个开黎曼面内的紧带边子流形，后文将看到这一点）。本小节的目标是回顾：当曲面
不必可定向时，可以在其上放置何种类型的结构——它容许一个一维复双解析流形结构，也称为
\emph{克莱因曲面}。克莱因曲面一般理论的标准参考书是 Alling 与 Greenleaf 的专著 [1]。
双解析结构近来在开 Gromov--Witten 理论中找到了应用，见 [10]。下文我们把「全纯」与
「解析」当作同义词使用。回忆：若 $U$ 是 $\C$ 的开子集，函数 $f:U\to\C$ 称为在 $U$
中\emph{解析}的，若 $\bar\partial f=0$；称为\emph{反解析}的，若 $\partial f=0$。

\begin{definition}[双解析函数，开情形]\label{def:2.1.9}
设 $U$ 是 $\C$ 的开子集。函数 $f:U\to\C$ 称为在 $U$ 中\emph{双解析}的，若它在 $U$
的任一给定连通分支 $V$ 上的限制 $f|_V$ 满足 $\bar\partial(f|_V)=0$ 或
$\partial(f|_V)=0$，即 $f$ 限制在该连通分支上要么解析、要么反解析。
\end{definition}

为了在本小节所考虑的一般性下的曲面上放置转移映射解析或双解析的图册，有必要把
双解析性推广到边界情形。

\begin{definition}[解析函数与双解析函数，边界情形]\label{def:2.1.10}
设 $A$ 是闭上半平面 $\C^+:=\{z\in\C\mid\operatorname{Im}z\ge0\}$ 的开子集。函数
$f:A\to\C$ 称为在 $A$ 上\emph{解析}（相应地，\emph{双解析}）的，若存在 $\C$ 的开子集
$U$ 使 $U\supset A$，且 $f$ 是某个 $U$ 上解析（相应地，双解析）函数在 $A$ 上的限制。
\end{definition}

下一个结果在联系克莱因曲面（定义 2.1.12）与 Real 黎曼面（定义 2.1.1）时是关键。

\begin{theorem}[Schwarz 反射原理，参看 [3]]\label{thm:2.1.11}
设 $A$ 是 $\C^+=\{z\in\C\mid\operatorname{Im}z\ge0\}$ 的开子集，$f:A\to\C$ 是满足
下列条件的函数：
\begin{enumerate}[label=(\roman*), leftmargin=2.2em, itemsep=0.15em]
\item $f$ 在 $A$ 上连续；
\item $f$ 在 $A$ 的内部解析；
\item $f(A\cap\R)\subset\R$。
\end{enumerate}
则存在唯一的解析函数 $F:A\cup\bar A\to\C$（定义在 $\C$ 的开子集 $A\cup\bar A$ 上），
使得
\begin{enumerate}[label=(\roman*), resume, leftmargin=2.2em, itemsep=0.15em]
\item $F|_A=f$；
\item 对一切 $z\in A\cup\bar A$，$F(\bar z)=\overline{F(z)}$。
\end{enumerate}
\end{theorem}

现在可以把克莱因曲面定义为自然推广黎曼面的对象。特别地，这将处理带边黎曼面的情形。

\begin{definition}[克莱因曲面]\label{def:2.1.12}
\emph{克莱因曲面}是一个曲面 $\Sigma$，带有一个图册 $(U_i,\phi_i)$，满足
\begin{enumerate}[label=(\roman*), leftmargin=2.2em, itemsep=0.15em]
\item $\phi_i:U_i\to V_i$ 是到 $\C^+=\{z\in\C\mid\operatorname{Im}z\ge0\}$ 的某个开
      子集 $V_i$ 上的同胚；
\item 若 $U_i\cap U_j\neq\varnothing$，则映射
      $\phi_i\circ\phi_j^{-1}:\phi_j(U_i\cap U_j)\to\phi_i(U_i\cap U_j)$ 是双解析的。
\end{enumerate}
这样的图册称为\emph{双解析图册}。
\end{definition}

\emph{带边黎曼面}是一个曲面 $\Sigma$，带有一个图册 $(U_i,\phi_i)$，满足
\begin{enumerate}[label=(\roman*), start=3, leftmargin=2.2em, itemsep=0.15em]
\item $\phi_i:U_i\to V_i$ 是到 $\C^+=\{z\in\C\mid\operatorname{Im}z\ge0\}$ 的某个开
      子集 $V_i$ 上的同胚；
\item 若 $U_i\cap U_j\neq\varnothing$，则映射
      $\phi_i\circ\phi_j^{-1}:\phi_j(U_i\cap U_j)\to\phi_i(U_i\cap U_j)$ 是解析的。
\end{enumerate}
这样的图册称为\emph{解析图册}。它只能存在于可定向曲面上。

于是，通常意义的黎曼面是带边黎曼面，而带边黎曼面是图册实为解析的克莱因曲面。作为
例子，$\C^+$ 的任一开子集 $A$ 是（可能带空边界的）带边黎曼面。更有趣的例子由下述
结果提供。

\begin{theorem}[Alling 与 Greenleaf, [1]]\label{thm:2.1.13}
设 $S$ 是紧连通拓扑曲面，可以不可定向、可以带边界。则 $S$ 上存在双解析图册。
若 $S$ 可定向，则其上存在解析图册。
\end{theorem}
